\begin{table}[H]
\scriptsize
\setlength{\tabcolsep}{3pt}
\renewcommand{\arraystretch}{0.95}

\caption{\label{tab:ea20-counterfactual-income-gap-product-contributions}Household price policies: product-group contributions to no-policy inflation inequality in the 20-country euro area, June 2021--June 2023}
\centering
\resizebox{\textwidth}{!}{%
\begin{tabular}{lrrrrr}
\toprule
Product group & \shortstack{Counterfactual\\ inflation} & \shortstack{Q1 counterfactual\\ contribution} & \shortstack{Q5 counterfactual\\ contribution} & \shortstack{Counterfactual\\ inflation gap} & \shortstack{Effect of price policies\\ on inflation inequality}\\
\midrule
Food and non-alcoholic beverages & 4.2 & 5.0 & 4.3 & 0.7 & 0.0\\
Alcoholic beverage, tobacco and narcotics & 0.4 & 0.6 & 0.4 & 0.2 & 0.0\\
Clothing and footwear & 0.3 & 0.3 & 0.4 & -0.2 & 0.0\\
Housing (rentals and repairs) & 0.4 & 0.4 & 0.6 & -0.2 & -0.0\\
Water, electricity, gas and other fuels & 7.4 & 10.1 & 7.1 & 3.0 & -1.5\\
\addlinespace
Transport & 2.3 & 1.7 & 3.1 & -1.3 & 0.2\\
Other & 0.3 & 0.2 & 0.4 & -0.1 & 0.0\\
\midrule
\textbf{Total} & \textbf{15.4} & \textbf{18.4} & \textbf{16.3} & \textbf{2.1} & \textbf{-1.2}\\
\bottomrule
\end{tabular}%
}
\vspace{-0.4em}
\begin{flushleft}
\footnotesize{Notes. The no-policy counterfactual replaces prices affected by household price policies and retains other observed prices. The gap is Q1 minus Q5. Policy effects are observed minus no-policy values; negative values indicate that policies reduced the gap. The euro-area aggregate covers all 20 countries using annual HICP country weights. Entries are percentage-point contributions. Totals are calculated before rounding and may differ from the sums of displayed entries. Sources: Eurostat, national statistical institutes, authors' calculations.}
\end{flushleft}
\end{table}
